\section{Тестирование программного обеспечения}

В данном разделе рассматривается тестирование программного обеспечения, разработанного для оценки кредитоспособности физических лиц с использованием алгоритма градиентного бустинга деревьев решений.

Целью тестирования являлась проверка корректности работы основных компонентов программного комплекса: ввода данных заемщика, формирования входного JSON-файла, запуска Python-модуля предсказания, получения результата скоринга и отображения результата в пользовательском интерфейсе.

Тестирование выполнялось на синтетических наборах данных, подготовленных для проверки различных сценариев работы приложения. Дополнительно проверялась возможность загрузки данных заемщика из JSON-файла.

\subsection{Функциональное тестирование}

Функциональное тестирование проводилось для проверки соответствия разработанного программного обеспечения заданным требованиям. В ходе тестирования проверялась работа основных функций системы, связанных с вводом данных заемщика, обработкой информации, запуском модуля предсказания и отображением результата кредитного скоринга.

При проведении функционального тестирования были рассмотрены следующие варианты работы программы:

\begin{itemize}
  \item ручной ввод данных заемщика через формы приложения;
  \item формирование JSON-структуры на основе введенных данных;
  \item загрузка данных заемщика из JSON-файла;
  \item запуск Python-модуля предсказания из C\#-приложения;
  \item получение результата работы модели в формате JSON;
  \item отображение вероятности дефолта, скорингового балла и уровня риска;
  \item отображение факторов риска и рекомендаций;
  \item обработка ошибок при отсутствии файла или некорректных входных данных.
\end{itemize}

Разработанное приложение принимает входные данные о заемщике, передает их в модуль предсказания, реализованный на языке Python, и получает результат оценки кредитоспособности. Взаимодействие между приложением и Python-модулем выполняется посредством запуска Python-скрипта через класс \texttt{ProcessStartInfo} с последующим чтением результата из стандартного потока вывода.

\subsubsection{Алгоритм тестирования}

Алгоритм функционального тестирования программного обеспечения на формальном языке можно представить следующим образом:

\begin{enumerate}
  \item Начало тестирования.
  \item Подготовить набор тестовых данных заемщика.
  \item Выбрать способ ввода данных:
        \begin{itemize}
          \item если используется ручной ввод, открыть формы приложения и заполнить поля заемщика;
          \item если используется загрузка из файла, выбрать JSON-файл с данными заемщика.
        \end{itemize}
  \item Проверить корректность заполнения обязательных полей.
  \item Если обязательные поля не заполнены или данные имеют некорректный формат, вывести сообщение об ошибке и перейти к шагу 13.
  \item Сформировать входную JSON-структуру для Python-модуля предсказания.
  \item Запустить Python-скрипт предсказания из C\#-приложения.
  \item Передать сформированный JSON-файл в модуль предсказания.
  \item Получить выходной JSON с результатами скоринга.
  \item Проверить наличие в выходных данных следующих полей:
        \begin{itemize}
          \item вероятность дефолта;
          \item скоринговый балл;
          \item уровень кредитного риска;
          \item решение о выдаче кредита;
          \item факторы, влияющие на решение модели.
        \end{itemize}
  \item Отобразить результат в пользовательском интерфейсе.
  \item Сравнить фактический результат работы программы с ожидаемым результатом.
  \item Зафиксировать результат тестирования.
  \item Конец тестирования.
\end{enumerate}

\subsubsection{Классы эквивалентности}

Для проверки корректности работы программного обеспечения были выделены классы эквивалентности входных данных. Это позволило сократить количество тестовых случаев и при этом проверить основные варианты поведения системы.

\begin{table}[ht!]
  \centering
  \caption{Классы эквивалентности входных данных}
  \label{tab:equivalence_classes}
  \begin{tabular}{|p{0.25\linewidth}|p{0.45\linewidth}|p{0.2\linewidth}|}
    \hline
    Класс данных                     & Описание класса                                                                        & Ожидаемый результат                                          \\
    \hline
    Корректные данные заемщика       & Все обязательные поля заполнены, числовые значения находятся в допустимых пределах     & Данные принимаются системой                                  \\
    \hline
    Незаполненные обязательные поля  & Отсутствует одно или несколько обязательных значений, необходимых для расчета скоринга & Пользователю выводится сообщение об ошибке                   \\
    \hline
    Некорректный формат данных       & В числовые поля введены текстовые значения или данные имеют неверный формат            & Данные не передаются в модель, выводится сообщение об ошибке \\
    \hline
    Корректный JSON-файл             & JSON-файл имеет правильную структуру и содержит необходимые поля                       & Файл загружается, данные передаются в модуль предсказания    \\
    \hline
    Некорректный JSON-файл           & JSON-файл поврежден или не соответствует требуемой структуре                           & Пользователю выводится сообщение об ошибке                   \\
    \hline
    Заемщик с низким уровнем риска   & Доход достаточен, кредитная нагрузка низкая, просрочки отсутствуют                     & Программа формирует результат скоринга без ошибок            \\
    \hline
    Заемщик со средним уровнем риска & Присутствует умеренная кредитная нагрузка или отдельные факторы риска                  & Программа формирует результат скоринга без ошибок            \\
    \hline
    Заемщик с высоким уровнем риска  & Имеются просрочки, высокая кредитная нагрузка или большое количество отказов           & Программа формирует результат скоринга без ошибок            \\
    \hline
  \end{tabular}
\end{table}

В качестве тестовых наборов использовались данные заемщиков с различными характеристиками: заемщик с низким уровнем риска, заемщик со средней кредитной нагрузкой, заемщик с высоким уровнем риска, пенсионер, а также заемщик с большим количеством просрочек и отказов. Такой набор тестов позволил проверить корректность заполнения всех полей формы, передачу данных в модуль предсказания и формирование итогового результата скоринга.

\FloatBarrier

\subsubsection{Результаты функционального тестирования}

Результаты проведенного функционального тестирования представлены в таблице~\ref{tab:functional_testing_results}.

\begin{table}[ht!]
  \centering
  \caption{Результаты функционального тестирования}
  \label{tab:functional_testing_results}
  \begin{tabular}{|p{0.38\linewidth}|p{0.42\linewidth}|p{0.12\linewidth}|}
    \hline
    Проверяемый сценарий              & Ожидаемый результат                                                                                & Результат \\
    \hline
    Ручной ввод данных заемщика       & Данные принимаются системой и проходят проверку заполнения                                         & Успешно   \\
    \hline
    Формирование JSON по данным формы & Создается JSON-структура, соответствующая формату входных данных модели                            & Успешно   \\
    \hline
    Загрузка данных из JSON-файла     & Файл считывается, данные передаются в модуль предсказания                                          & Успешно   \\
    \hline
    Запуск Python-скрипта из C\#      & Python-модуль запускается из приложения и получает путь к входному файлу                           & Успешно   \\
    \hline
    Получение результата скоринга     & Приложение получает JSON с вероятностью дефолта, баллом и уровнем риска                            & Успешно   \\
    \hline
    Отображение результата            & На форме отображаются вероятность дефолта, скоринговый балл, уровень риска, факторы и рекомендации & Успешно   \\
    \hline
    Обработка некорректных данных     & Пользователю выводится сообщение об ошибке без аварийного завершения программы                     & Успешно   \\
    \hline
  \end{tabular}
\end{table}

\FloatBarrier

По результатам функционального тестирования установлено, что разработанное программное обеспечение корректно выполняет основные пользовательские сценарии. Приложение обеспечивает ручной ввод данных заемщика, загрузку данных из JSON-файла, формирование входной структуры данных, запуск Python-модуля предсказания, получение результата скоринга и отображение итоговой оценки пользователю.

При проверке некорректных входных данных программа не завершалась аварийно, а выводила сообщение об ошибке. Это подтверждает корректность обработки ошибочных ситуаций на уровне пользовательского сценария.

Функциональное тестирование подтвердило корректность работы основных компонентов программного комплекса.

\subsection{Модульное тестирование}

Модульное тестирование не проводилось.

\subsection{Тестирование графического интерфейса}

Тестирование графического интерфейса не проводилось.

\subsection*{Выводы}

В ходе тестирования была проверена корректность работы программного обеспечения для оценки кредитоспособности физических лиц. Были протестированы основные функциональные сценарии: ручной ввод данных заемщика, загрузка данных из JSON-файла, формирование входной структуры данных, запуск Python-модуля предсказания, получение результата скоринга и отображение результата в пользовательском интерфейсе.

Результаты функционального тестирования показали, что разработанное программное обеспечение выполняет основные функции корректно. Приложение принимает данные заемщика, формирует JSON-структуру, передает данные в Python-модуль и получает результат скоринга без нарушения общего сценария работы.

Проверена обработка ошибочных ситуаций, связанных некорректными входными данными или отсутствием файла. В таких случаях пользователю выводится сообщение об ошибке, а программа не завершается аварийно.

Модульное тестирование и тестирование графического интерфейса в рамках практики не проводились.